% !TEX program = lualatex
%!TEX TS-program = lualatex
\documentclass[12pt,a4paper]{article}
\usepackage{fontspec}
\usepackage{polyglossia}
\setdefaultlanguage{polish}
\usepackage{geometry}
\geometry{margin=2.1cm}
\usepackage[table]{xcolor}
\usepackage{longtable}
\usepackage{array}
\usepackage{booktabs}
\usepackage{enumitem}
\usepackage{ragged2e}
\definecolor{tableHeader}{HTML}{E8EEF7}
\definecolor{tableStripe}{HTML}{F7F9FC}
\newcolumntype{L}[1]{>{\RaggedRight\arraybackslash}p{#1}}
\newcolumntype{C}[1]{>{\Centering\arraybackslash}p{#1}}
\renewcommand{\arraystretch}{1.22}
\setlength{\tabcolsep}{4pt}
\setlist[itemize]{leftmargin=*, itemsep=2pt, topsep=2pt}
\setlength{\parindent}{0pt}
\setlength{\parskip}{6pt}
\emergencystretch=3em
\sloppy
\begin{document}

\section*{Raport ze sprintu nr 2 - Calisia Banking App}
Data raportu: 08.05.2026

Okres sprintu: 03.05 - 08.05.2026

Product Owner / Scrum Master: Jakub Łaszuk

\subsection*{Lista osób}
\begin{itemize}
\item Kubryn Jeremiasz
\item Łaszuk Jakub
\item Makarevich Ksenia
\item Nuckowski Michał
\item Ruczyńska Nikola
\item Wołczko Jakub
\item Olszewski Dawid
\end{itemize}

\subsection*{Product Owner i Scrum Master sprintu}
Jakub Łaszuk

\subsection*{Lista historyjek na sprint}
\begingroup
\footnotesize
\rowcolors{2}{tableStripe}{white}
\begin{longtable}{@{}C{0.08\textwidth}L{0.62\textwidth}C{0.10\textwidth}L{0.12\textwidth}@{}}
\toprule
\rowcolor{tableHeader}
\textbf{Kod} & \textbf{Historyjka} & \textbf{SP} & \textbf{Status} \\
\midrule
\endfirsthead
\toprule
\rowcolor{tableHeader}
\textbf{Kod} & \textbf{Historyjka} & \textbf{SP} & \textbf{Status} \\
\midrule
\endhead
S2-01 & Rejestracja profilu klienta z imieniem, nazwiskiem, e-mailem, hasłem i typem klienta & 8 & Done \\
S2-02 & Walidacja danych rejestracyjnych & 5 & Done \\
S2-03 & Bezpieczne haszowanie haseł z użyciem PBKDF2 & 5 & Done \\
S2-04 & Jednoznaczna odpowiedź po rejestracji klienta & 3 & Done \\
S2-05 & Testy jednostkowe i integracyjne rejestracji & 5 & Done \\
S2-06 & Scenariusz rejestracji w Bruno & 2 & Done \\
\bottomrule
\end{longtable}
\endgroup

\subsection*{Podsumowanie sprintu}
Cel sprintu: Klient, rejestracja i bezpieczne hasła. Umożliwienie klientowi założenia profilu oraz przygotowanie bezpiecznego przechowywania danych logowania.

Zakres sprintu obejmował 6 historyjek o łącznej estymacji 28 SP. Wszystkie historyjki zostały dowiezione ze statusem Done, dlatego osiągnięta prędkość sprintu wyniosła 28 SP.

Dostarczono przepływ rejestracji klienta przez API wraz z walidacją danych, obsługą typu klienta, bezpiecznym haszowaniem hasła PBKDF2, odpowiedzią zawierającą identyfikator klienta oraz testami i scenariuszem Bruno.

Sprint domknął pierwszy pełny proces biznesowy klienta. System zaczął przyjmować dane użytkownika w kontrolowany sposób, a mechanizm PBKDF2 ograniczył ryzyko związane z przechowywaniem danych logowania.

\subsubsection*{Najważniejsze rezultaty}
\begin{itemize}
\item Utworzono profil klienta przez API.
\item Dodano walidację danych rejestracyjnych i obsługę błędów.
\item Wdrożono bezpieczne haszowanie haseł.
\item Przygotowano testy oraz scenariusz rejestracji w Bruno.
\end{itemize}

\subsection*{Retrospekcja sprintu}
Retrospekcja potwierdziła, że zespół dobrze radzi sobie z funkcjami przekrojowymi, które dotykają domeny, API, bazy i testów.

\subsubsection*{Co się udało}
\begin{itemize}
\item Sprawna implementacja komendy rejestracji klienta zgodnie z CQRS.
\item Dobre pokrycie przypadków walidacyjnych testami.
\item Wprowadzenie PBKDF2 bez przechowywania haseł w postaci jawnej.
\item Bieżąca aktualizacja kolekcji Bruno.
\end{itemize}
\subsubsection*{Poziom energii zespołu}
Średni poziom energii zespołu: 4,2/5.

Energia była wysoka dzięki realistycznemu zakresowi sprintu i satysfakcji z pierwszego kompletnego procesu biznesowego.

\subsubsection*{Czego udało się nauczyć}
\begin{itemize}
\item Zespół przećwiczył bezpieczne przechowywanie haseł.
\item Warto wcześniej ustalać reguły walidacji danych klienta.
\item Kontrakty API powinny być synchronizowane z formularzami frontendu.
\end{itemize}
\subsubsection*{Wnioski i działania na kolejny sprint}
\begin{itemize}
\item Przygotować wspólną listę reguł walidacyjnych dla danych klienta.
\item Dodawać przykłady request/response do Bruno przy każdym nowym endpoincie.
\item Utrzymać testy negatywne dla procesów logowania i rejestracji.
\end{itemize}

\subsection*{Następny sprint}
Product Owner / Scrum Master w następnym sprincie: Jeremiasz Kubryn

Główny cel następnego sprintu: Logowanie, JWT i ochrona API. Zapewnienie mechanizmu uwierzytelniania oraz zabezpieczenie operacji klienta przed dostępem anonimowym.

\end{document}
